\section*{Overview}

For hard 64-bit factorization tasks, trial division is not enough. The standard contest toolkit is Miller-Rabin for
primality testing and Pollard Rho for splitting composite numbers into nontrivial factors.

\section*{When to Use It}

Use this pair when:

\begin{itemize}
  \item inputs are large 64-bit integers,
  \item repeated primality tests appear,
  \item full factorization is required and \(\sqrt{n}\) trial division is too slow.
\end{itemize}

\section*{Core Idea}

Miller-Rabin is a fast compositeness test based on modular exponentiation. Pollard Rho is a randomized splitting method
that tries to find a nontrivial gcd of values generated by a polynomial walk modulo \(n\).

Together, they form a recursive factorizer:

\begin{itemize}
  \item if \(n\) is prime by Miller-Rabin, stop,
  \item otherwise use Pollard Rho to find one factor \(d\),
  \item recurse on \(d\) and \(n/d\).
\end{itemize}

\section*{Key Insight}

Miller-Rabin answers ``is this number prime?'' quickly enough that Pollard Rho only spends real effort on actual
composites. Pollard Rho does not need to find the smallest factor; it only needs any nontrivial factor so recursion can
continue.

\section*{Operations / Main Technique}

\begin{itemize}
  \item safe modular multiplication,
  \item fast modular power,
  \item deterministic Miller-Rabin for 64-bit integers using a fixed witness set,
  \item randomized Pollard Rho splitting,
  \item recursive factor collection.
\end{itemize}

\paragraph{Worked example.}
If \(n = pq\) with two large odd primes, Miller-Rabin quickly says ``composite.'' Pollard Rho then runs a modular walk
until the gcd of two iterates reveals either \(p\) or \(q\).

\section*{Correctness Intuition}

Miller-Rabin works because primes satisfy strong modular identities that composites often violate. Pollard Rho works by
forcing two sequences modulo an unknown factor to collide earlier than they do modulo the full composite, so the gcd of
their difference with \(n\) leaks that factor.

\section*{Complexity Analysis}

The exact running time of Pollard Rho is probabilistic, but in practice it is fast enough for standard 64-bit contest
inputs. Miller-Rabin itself is \(O(k \log n)\) modular multiplications for a fixed witness count.

\section*{Implementation}

The code handles 64-bit integers with:

\begin{itemize}
  \item \texttt{\_\_int128} modular arithmetic,
  \item deterministic 64-bit Miller-Rabin witnesses,
  \item recursive factorization that returns the prime factors sorted.
\end{itemize}

\section*{Common Pitfalls}

\begin{itemize}
  \item Overflow in modular multiplication without \texttt{\_\_int128}.
  \item Using too weak a witness set for 64-bit determinism.
  \item Forgetting to handle even numbers and small primes separately.
  \item Treating Pollard Rho as deterministic; it may need a restart with a new constant.
\end{itemize}

\section*{Variants / Extensions}

\begin{itemize}
  \item Probabilistic Miller-Rabin for bigger integer types.
  \item Use Pollard Rho only on the leftovers after small-prime trial division.
  \item Build divisor generation or multiplicative functions on top of the returned prime factors.
\end{itemize}

\section*{Practice Problems}

\begin{itemize}
  \item Factor arbitrary 64-bit integers.
  \item Answer many primality queries on large numbers.
  \item Combine factorization with divisor or multiplicative-function tasks.
\end{itemize}

\section*{References}

\begin{itemize}
  \item \href{https://cp-algorithms.com/algebra/factorization.html}{cp-algorithms: Integer Factorization}.
  \item \href{https://codeforces.com/catalog?locale=en}{Codeforces Catalog}.
  \item \href{https://github.com/kth-competitive-programming/kactl}{KACTL}.
\end{itemize}
